\section{Related Work}
\label{sec:related}

\paragraph{Chunk enrichment.}
Retrieval quality is decided as much by what is indexed as by the encoder.
Document expansion by query prediction appends generated queries to a passage
before indexing \citep{nogueira2019document}, and proposition-level indexing
replaces passages with self-contained statements
\citep{chen2024densex}. A recent taxonomy and evaluation of chunking strategies
shows that the choice of chunking policy interacts strongly with the retriever
\citep{zhou2026chunking}, while late chunking embeds a whole document before
splitting it, so that every chunk vector already carries document context
\citep{mortezaagha2026latechunking}. Contextual retrieval works along the same
lines, with a document-conditioned summary prepended to each chunk
\citep{anthropic2024contextual}. These methods are usually evaluated together
with a cross-encoder reranker, and the published recipe reports the two
components jointly, which makes their individual contributions hard to
disentangle. We hold first-stage retrieval fixed and vary only the index text and the
presence of the reranker, and we go one step further by distilling the
enrichment model itself so that the ingest stage also runs on-device.

\paragraph{Reranking and late interaction.}
Cross-encoder reranking is an established second stage over a first-stage
candidate list \citep{nogueira2019passage}, and late-interaction models such as
ColBERT retain token-level matching without a full cross-encoder
\citep{khattab2020colbert}. A controlled study of long-context question
answering finds that learned context planning and reranking do not improve over
strong first-stage retrieval \citep{li2026reranking}, which agrees with our
negative result. Modern instruction-tuned embedders and rerankers supply both
stages from a single model family \citep{qwen3embedding}. Reranking is usually
reported as effective, but it scales linearly in the candidate pool and is the
most expensive part of an interactive pipeline on a laptop. We quantify that
cost on a fixed device, and find that once the chunks are contextualized the
reranker can be dropped with no measurable loss in quality.

\paragraph{Hybrid retrieval and fusion.}
Sparse lexical retrieval with BM25 remains a strong baseline
\citep{robertson2009bm25}, and learned sparse models such as SPLADE improve on
it \citep{formal2021splade}. Dense and sparse lists are commonly combined with
reciprocal rank fusion \citep{cormack2009reciprocal}, and self-contained
tooling for hybrid retrieval over quantized embeddings and lexical indexes is
available \citep{livieris2026quanta}. Hybrid retrieval is reported to help on
heterogeneous corpora \citep{luan2021sparse}, but the benefit depends on how
much the two channels complement each other. Our corpus is homogeneous, and we
find that fusion hurts consistently, including with a stronger sparse backend.

\paragraph{Distillation and parameter-efficient tuning.}
Knowledge distillation transfers the behavior of a large model to a smaller
one \citep{hinton2015distilling}, and low-rank adaptation makes the transfer
cheap \citep{hu2022lora,dettmers2023qlora}. We use LoRA over a 4-bit base model
to distill an API teacher into 0.6B parameters, and we judge the student by the
retrieval quality it supports rather than by generation quality alone.
Distilling a retrieval-augmented pipeline for smartphone deployment has been
studied in the legal domain \citep{kamal2026smartphone}, and on-premise
deployments select sub-networks that fit the local hardware
\citep{rizeakos2026onprem}.

\paragraph{Efficient retrieval on device.}
Matryoshka representation learning trains embeddings whose prefixes are
individually usable, which allows truncation with little loss
\citep{kusupati2022matryoshka}; dimension-elastic indexes serve several
truncation levels from one resident index \citep{culatana2026mrvq}, and
sentence encoders provide the dense backbone
\citep{reimers2019sentence,karpukhin2020dense}. Retrieval-augmented generation
\citep{lewis2020retrieval} and its surveys \citep{gao2023retrieval} describe
the full pipeline; the practical constraint in our setting is that every
component must fit and run on one consumer device. We therefore report latency
and cost alongside quality, and we build on the MLX framework for Apple silicon
\citep{mlx} with a local vector index \citep{qdrant}.

